% ---- score2 standard preamble (from: python3 score2_tools.py preamble) ----
\documentclass[10pt]{extarticle}
\usepackage[a4paper,margin=1.5cm]{geometry}
\usepackage[T1]{fontenc}
\usepackage{amsmath,amssymb}
\usepackage{graphicx}
\usepackage[export]{adjustbox}
\usepackage{multicol}
\usepackage{enumitem}
\usepackage{booktabs,array,tabularx}
\setlength{\columnsep}{0.7cm}
\setlength{\parindent}{0pt}
\setlength{\parskip}{3pt}
\pagestyle{empty}
\makeatletter
% figure registry (figdefs.tex is generated by the extract step)
\newcommand{\qfigdef}[4]{\@namedef{qf@file@#1}{#2}\@namedef{qf@w@#1}{#3}\@namedef{qf@wide@#1}{#4}}
\newcommand{\qf@check}[1]{\@ifundefined{qf@file@#1}{\PackageError{score2}{Unknown figure id #1}{}}{}}
\newcommand{\qf@img}[4][]{\includegraphics[width=#2\dimexpr\@nameuse{qf@w@#4}pt\relax,max width=#3,max height=0.45\textheight,#1]{\@nameuse{qf@file@#4}}}
% Full-width questions: finalize wraps every question that contains a wide figure in \qfullbegin...\qfullend,
% so the WHOLE question (stem, figures, options) is set across the page in reading order, then the two columns
% resume. The question's head - its number, stem and first figure (\qfullsplit marks the end of the head) -
% is kept on one page, so a number is never left at a page foot with its figure overleaf; the rest may break
% across pages normally, so no big gaps are left. Without a \qfullsplit, the whole question is kept together
% if it is under 45% of a page. Never write these yourself.
\newif\ifqf@full
\newif\ifqf@boxed
\newsavebox{\qf@box}
\newcommand{\qf@open}{\global\qf@fulltrue\par\global\qf@boxedtrue
  \setbox\qf@box\vbox\bgroup\hsize\textwidth\linewidth\textwidth}
\newcommand{\qf@emit}[1]{\par\egroup\global\qf@boxedfalse
  \ifdim\dimexpr\ht\qf@box+\dp\qf@box\relax<#1\textheight\noindent\usebox{\qf@box}\par\else\unvbox\qf@box\fi}
\newcommand{\qfullsplit}{\ifqf@boxed\qf@emit{0.85}\fi}
\newcommand{\qf@close}{\ifqf@boxed\qf@emit{0.45}\fi\global\qf@fullfalse}
\newcommand{\qfullbegin}{\end{multicols}\qf@open}
\newcommand{\qfullend}{\qf@close\begin{multicols}{2}}
\newcommand{\qfullstart}{\qf@open}
\newcommand{\qfulldone}{\qf@close}
% \qfig[scale]{id}: block figure, centred, directly where it is written (size set by the tool).
% Fallback if a wide figure is not inside a full-width question: break out of the columns in place - no floats.
\newcommand{\qfig}[2][1]{\qf@check{#2}%
  \ifqf@full\qf@blk{#1}{#2}\else
  \ifnum\@nameuse{qf@wide@#2}=1 \expandafter\@firstoftwo\else\expandafter\@secondoftwo\fi
  {\end{multicols}\begin{center}\qf@img{#1}{\textwidth}{#2}\end{center}\begin{multicols}{2}}%
  {\qf@blk{#1}{#2}}\fi}
\newcommand{\qf@blk}[2]{\par\smallskip{\centering\qf@img{#1}{\linewidth}{#2}\par}\smallskip}
% \qfigi[scale]{id}: inline figure (inside option lists / tables / minipages); never breaks columns.
\newcommand{\qfigi}[2][1]{\qf@check{#2}\qf@img[valign=c]{#1}{\linewidth}{#2}}
\makeatother
% structure macros - lint reads these to check completeness and figure placement
\newcommand{\papersec}[1]{\par\bigskip{\centering\Large\bfseries PAPER-#1\par}\medskip}
\newcommand{\partsec}[2]{\par\medskip\noindent{\large\bfseries PART-#1: #2}\par\smallskip}
\newcommand{\sect}[1]{\par\smallskip\noindent\textbf{#1}\par}
\newcommand{\qn}[1]{\par\smallskip\noindent\textbf{#1.}~}
\newcommand{\paraq}[2]{\par\smallskip\noindent\textbf{Paragraph for Questions #1 and #2}\par}
\newenvironment{qoptions}{\begin{itemize}[leftmargin=1.8em,itemsep=1pt,topsep=2pt,label={}]}{\end{itemize}}
\newcommand{\opt}[2]{\item[(#1)] #2}
% figure registry (generated by score2_tools.py finalize)
\qfigdef{p1_phy_s1_q01_1}{TOAS - Test 08 (15 May 2022)-img/p1_phy_s1_q01_1.png}{50.4}{0}
\qfigdef{p1_phy_s1_q02_1}{TOAS - Test 08 (15 May 2022)-img/p1_phy_s1_q02_1.png}{146.3}{0}
\qfigdef{p1_phy_s1_q03_1}{TOAS - Test 08 (15 May 2022)-img/p1_phy_s1_q03_1.png}{122.0}{0}
\qfigdef{p1_phy_s1_q04_1}{TOAS - Test 08 (15 May 2022)-img/p1_phy_s1_q04_1.png}{155.5}{0}
\qfigdef{p1_phy_s1_q05_A}{TOAS - Test 08 (15 May 2022)-img/p1_phy_s1_q05_A.png}{103.8}{0}
\qfigdef{p1_phy_s1_q05_B}{TOAS - Test 08 (15 May 2022)-img/p1_phy_s1_q05_B.png}{94.5}{0}
\qfigdef{p1_phy_s1_q05_C}{TOAS - Test 08 (15 May 2022)-img/p1_phy_s1_q05_C.png}{79.8}{0}
\qfigdef{p1_phy_s1_q05_D}{TOAS - Test 08 (15 May 2022)-img/p1_phy_s1_q05_D.png}{83.3}{0}
\qfigdef{p1_phy_s1_q06_1}{TOAS - Test 08 (15 May 2022)-img/p1_phy_s1_q06_1.png}{106.5}{0}
\qfigdef{p1_phy_s2_para7_8_1}{TOAS - Test 08 (15 May 2022)-img/p1_phy_s2_para7_8_1.png}{131.5}{0}
\qfigdef{p1_phy_s2_para9_10_1}{TOAS - Test 08 (15 May 2022)-img/p1_phy_s2_para9_10_1.png}{112.3}{0}
\qfigdef{p1_phy_s3_q01_1}{TOAS - Test 08 (15 May 2022)-img/p1_phy_s3_q01_1.png}{106.6}{0}
\qfigdef{p1_phy_s3_q03_1}{TOAS - Test 08 (15 May 2022)-img/p1_phy_s3_q03_1.png}{185.4}{0}
\qfigdef{p1_phy_s3_q05_1}{TOAS - Test 08 (15 May 2022)-img/p1_phy_s3_q05_1.png}{154.2}{0}
\qfigdef{p1_phy_s3_q06_1}{TOAS - Test 08 (15 May 2022)-img/p1_phy_s3_q06_1.png}{129.3}{0}
\qfigdef{p1_phy_s3_q07_1}{TOAS - Test 08 (15 May 2022)-img/p1_phy_s3_q07_1.png}{183.5}{0}
\qfigdef{p1_phy_s3_q08_1}{TOAS - Test 08 (15 May 2022)-img/p1_phy_s3_q08_1.png}{114.6}{0}
\qfigdef{p1_chem_s1_q04_1}{TOAS - Test 08 (15 May 2022)-img/p1_chem_s1_q04_1.png}{117.8}{0}
\qfigdef{p1_chem_s1_q05_A}{TOAS - Test 08 (15 May 2022)-img/p1_chem_s1_q05_A.png}{101.4}{0}
\qfigdef{p1_chem_s1_q05_B}{TOAS - Test 08 (15 May 2022)-img/p1_chem_s1_q05_B.png}{134.7}{0}
\qfigdef{p1_chem_s1_q06_A}{TOAS - Test 08 (15 May 2022)-img/p1_chem_s1_q06_A.png}{114.8}{0}
\qfigdef{p1_chem_s1_q06_B}{TOAS - Test 08 (15 May 2022)-img/p1_chem_s1_q06_B.png}{115.3}{0}
\qfigdef{p1_chem_s1_q06_C}{TOAS - Test 08 (15 May 2022)-img/p1_chem_s1_q06_C.png}{160.6}{0}
\qfigdef{p1_chem_s1_q06_D}{TOAS - Test 08 (15 May 2022)-img/p1_chem_s1_q06_D.png}{157.3}{0}
\qfigdef{p1_chem_s2_para9_10_1}{TOAS - Test 08 (15 May 2022)-img/p1_chem_s2_para9_10_1.png}{159.0}{0}
\qfigdef{p1_chem_s2_q10_A}{TOAS - Test 08 (15 May 2022)-img/p1_chem_s2_q10_A.png}{76.4}{0}
\qfigdef{p1_chem_s2_q10_B}{TOAS - Test 08 (15 May 2022)-img/p1_chem_s2_q10_B.png}{61.8}{0}
\qfigdef{p1_chem_s2_q10_C}{TOAS - Test 08 (15 May 2022)-img/p1_chem_s2_q10_C.png}{59.3}{0}
\qfigdef{p1_chem_s3_q03_1}{TOAS - Test 08 (15 May 2022)-img/p1_chem_s3_q03_1.png}{103.8}{0}
\qfigdef{p1_chem_s3_q07_1}{TOAS - Test 08 (15 May 2022)-img/p1_chem_s3_q07_1.png}{148.7}{0}
\qfigdef{p1_chem_s3_q07_2}{TOAS - Test 08 (15 May 2022)-img/p1_chem_s3_q07_2.png}{115.3}{0}
\qfigdef{p1_math_s3_q07_1}{TOAS - Test 08 (15 May 2022)-img/p1_math_s3_q07_1.png}{106.8}{0}
\qfigdef{p2_phy_s1_q01_1}{TOAS - Test 08 (15 May 2022)-img/p2_phy_s1_q01_1.png}{149.2}{0}
\qfigdef{p2_phy_s1_q02_1}{TOAS - Test 08 (15 May 2022)-img/p2_phy_s1_q02_1.png}{158.0}{0}
\qfigdef{p2_phy_s1_q03_1}{TOAS - Test 08 (15 May 2022)-img/p2_phy_s1_q03_1.png}{104.1}{0}
\qfigdef{p2_phy_s1_q04_1}{TOAS - Test 08 (15 May 2022)-img/p2_phy_s1_q04_1.png}{102.4}{0}
\qfigdef{p2_phy_s1_q05_1}{TOAS - Test 08 (15 May 2022)-img/p2_phy_s1_q05_1.png}{133.2}{0}
\qfigdef{p2_phy_s1_q06_1}{TOAS - Test 08 (15 May 2022)-img/p2_phy_s1_q06_1.png}{101.4}{0}
\qfigdef{p2_phy_s2_q07_1}{TOAS - Test 08 (15 May 2022)-img/p2_phy_s2_q07_1.png}{63.7}{0}
\qfigdef{p2_phy_s2_q08_1}{TOAS - Test 08 (15 May 2022)-img/p2_phy_s2_q08_1.png}{91.6}{0}
\qfigdef{p2_phy_s2_q09_1}{TOAS - Test 08 (15 May 2022)-img/p2_phy_s2_q09_1.png}{183.7}{0}
\qfigdef{p2_phy_s3_q01_1}{TOAS - Test 08 (15 May 2022)-img/p2_phy_s3_q01_1.png}{114.4}{0}
\qfigdef{p2_phy_s3_q03_1}{TOAS - Test 08 (15 May 2022)-img/p2_phy_s3_q03_1.png}{81.8}{0}
\qfigdef{p2_phy_s3_q04_1}{TOAS - Test 08 (15 May 2022)-img/p2_phy_s3_q04_1.png}{76.2}{0}
\qfigdef{p2_phy_s3_q05_1}{TOAS - Test 08 (15 May 2022)-img/p2_phy_s3_q05_1.png}{127.7}{0}
\qfigdef{p2_phy_s3_q08_1}{TOAS - Test 08 (15 May 2022)-img/p2_phy_s3_q08_1.png}{99.7}{0}
\qfigdef{p2_chem_s1_q05_A}{TOAS - Test 08 (15 May 2022)-img/p2_chem_s1_q05_A.png}{113.1}{0}
\qfigdef{p2_chem_s1_q05_B}{TOAS - Test 08 (15 May 2022)-img/p2_chem_s1_q05_B.png}{134.5}{0}
\qfigdef{p2_chem_s1_q05_D}{TOAS - Test 08 (15 May 2022)-img/p2_chem_s1_q05_D.png}{111.0}{0}
\qfigdef{p2_chem_s1_q06_1}{TOAS - Test 08 (15 May 2022)-img/p2_chem_s1_q06_1.png}{184.7}{0}
\qfigdef{p2_chem_s2_q10_1}{TOAS - Test 08 (15 May 2022)-img/p2_chem_s2_q10_1.png}{132.0}{0}
\qfigdef{p2_chem_s2_q10_2}{TOAS - Test 08 (15 May 2022)-img/p2_chem_s2_q10_2.png}{61.8}{0}
\qfigdef{p2_chem_s2_q10_3}{TOAS - Test 08 (15 May 2022)-img/p2_chem_s2_q10_3.png}{76.9}{0}
\qfigdef{p2_chem_s2_q10_4}{TOAS - Test 08 (15 May 2022)-img/p2_chem_s2_q10_4.png}{104.3}{0}
\qfigdef{p2_chem_s3_q06_1}{TOAS - Test 08 (15 May 2022)-img/p2_chem_s3_q06_1.png}{245.2}{0}
\qfigdef{p2_chem_s3_q07_1}{TOAS - Test 08 (15 May 2022)-img/p2_chem_s3_q07_1.png}{144.0}{0}
\qfigdef{p2_chem_s3_q08_1}{TOAS - Test 08 (15 May 2022)-img/p2_chem_s3_q08_1.png}{98.2}{0}
\raggedcolumns
\makeatletter\@beginparpenalty=10000\makeatother
\newcolumntype{L}[1]{>{\raggedright\arraybackslash}p{#1}}
\begin{document}
\begin{center}{\Large\bfseries ALLEN TOAS --- Test 08 (15 May 2022)}\end{center}
\begin{multicols}{2}
\raggedcolumns
\papersec{1}
\partsec{1}{PHYSICS}
\sect{SECTION-I(i)}
\qn{1} Two identical balls each of mass $m$ carries equal charge $Q$ are connected to a massless spring of spring constant $k$ and released from rest. Initially spring was in it natural length $\ell$. Take $k=\dfrac{2Q^2}{3\pi\varepsilon_0\ell^3}$, then select correct option(s)
\qfig{p1_phy_s1_q01_1}
\begin{qoptions}
\opt{A}{Maximum distance between balls in motion will be $\dfrac{3\ell}{2}$.}
\opt{B}{Both ball will perform oscillation in frame of centre of mass but these oscillation will not be simple harmonic.}
\opt{C}{Acceleration of centre of mass of system will be decreasing with time.}
\opt{D}{Acceleration of centre of mass of system will be increasing with time.}
\end{qoptions}
\qn{2} All the surface are smooth and collisions are elastic in nature. The system is released from rest. Then
\qfig{p1_phy_s1_q02_1}
\begin{qoptions}
\opt{A}{When block first time slide down wedge and reaches horizontal surface velocity of centre of mass of block-wedge system have zero velocity.}
\opt{B}{After a long time for $n<3$ velocity of centre of mass of block wedge system will be $\sqrt{\dfrac{8ghn}{(n+1)^3}}$.}
\opt{C}{After a long time for $4>n>3$ velocity of centre of mass of block wedge system will be $\sqrt{\dfrac{32ghn(n-1)^2}{(n+1)^5}}$.}
\opt{D}{Velocity of centre of mass of wedge system will not have vertical component.}
\end{qoptions}
\qn{3} A point source $S$ is placed in water at origin $(0,0)$ and an observer $O$ is in air at $(7\text{m},7\text{m})$ as shown in figure. (Given : $\mu_{\text{Water}}=4/3$)
\qfig{p1_phy_s1_q03_1}
\begin{qoptions}
\opt{A}{Ray starting from source along $3\hat i+4\hat j$ vector will reach observer after refraction.}
\opt{B}{Image of object seen by observer will be at line $4y-3x=7$.}
\opt{C}{Co-ordinate of image seen by observer is $\left(\dfrac{21}{16}\text{m},\dfrac{175}{64}\text{m}\right)$.}
\opt{D}{Ray starting from source along $4\hat i+3\hat j$ vector will reach observer after refraction.}
\end{qoptions}
\qn{4} In given circuit all batteries are ideal than:
\qfig{p1_phy_s1_q04_1}
\begin{qoptions}
\opt{A}{Current in branch CD is 1.4 A from D to C.}
\opt{B}{10V battery is charging, consuming 6W power.}
\opt{C}{Rate of heat production in branch CH is 8.4 J/s.}
\opt{D}{In this circuit only 20V battery is discharging supplying power of 98W.}
\end{qoptions}
\qn{5} Assuming the earth to be uniform sphere of radius $R$ and mass $M$ and neglecting rotational motion select correct option(s)
\begin{qoptions}
\opt{A}{\qfigi{p1_phy_s1_q05_A}\newline $|\vec g|$ : Acceleration due to gravity, $r$ : distance from centre}
\opt{B}{\qfigi{p1_phy_s1_q05_B}\newline $v$ : gravitational potential}
\opt{C}{\qfigi{p1_phy_s1_q05_C}\newline $P$ : gravitational pressure.}
\opt{D}{\qfigi{p1_phy_s1_q05_D}\newline $T$ : time period of simple pendulum.}
\end{qoptions}
\columnbreak
\qn{6} Three conducting spherical shell $A$,$B$ and $C$ has radius $R$, $2R$ and $3R$ respectively. Conductor $B$ is connected to earth and $A$ and $C$ are connected by conducting wire. If charge $Q$ is given to conductor $C$ then :-
\qfig{p1_phy_s1_q06_1}
\begin{qoptions}
\opt{A}{Charge on outer surface of conductor $A$ is $\dfrac{2Q}{11}$.}
\opt{B}{Charge on inner surface of conductor $C$ is $\dfrac{6Q}{11}$.}
\opt{C}{Charge on conductor $B$ is $-\dfrac{8Q}{11}$.}
\opt{D}{Charge on outer surface of conductor $C$ is $\dfrac{3Q}{11}$}
\end{qoptions}
\sect{SECTION-I(ii)}
\paraq{7}{8}
A Non-conducting ring of mass $m$ and radius $R$ is placed on a rough horizontal surface as shown in figure.
\qfig{p1_phy_s2_para7_8_1}
It has linear charge density $\lambda=\lambda_0\sin\phi$ where $\lambda_0$ is constant and angle $\phi$ is angle of radius vector with line $OA$.
\qn{7} Electric field intensity at centre of ring is :
\begin{qoptions}
\opt{A}{$\dfrac{\lambda_0}{8\pi\varepsilon_0R}(\hat i)$}
\opt{B}{$\dfrac{\lambda_0}{4\varepsilon_0R}(-\hat j)$}
\opt{C}{$\dfrac{\lambda_0}{8\varepsilon_0R}(-\hat j)$}
\opt{D}{$\dfrac{\lambda_0}{4\pi^2\varepsilon_0R}(+\hat j)$}
\end{qoptions}
\qn{8} When external uniform electric field $\vec E=E_0\hat i$ is switch on at $t=0$ then speed of ring after it rotates by angel $\pi/2$ will be (assuming friction is sufficient to prevent slipping]
\begin{qoptions}
\opt{A}{$R\sqrt{\dfrac{3\pi\lambda_0E_0}{m}}$}
\opt{B}{$R\sqrt{\dfrac{\lambda_0E_0}{m}}$}
\opt{C}{$R\sqrt{\dfrac{\pi\lambda_0E_0}{m}}$}
\opt{D}{$\pi R\sqrt{\dfrac{2\lambda_0E_0}{m}}$}
\end{qoptions}
\paraq{9}{10}
A thick plano-convex lens is made of material having refractive index $\mu=\dfrac32$. When it's plane surface is slivered and an object is placed 60 cm away from it's convex surface then final image forms at object it self. While when it's convex surface is silvered and object is placed at distance $\dfrac{20}{3}$cm away from plane surface then final image again forms at object itself.
\qfig{p1_phy_s2_para9_10_1}
\qn{9} Radius of curvature of convex surface is :
\begin{qoptions}
\opt{A}{10 cm}
\opt{B}{15 cm}
\opt{C}{25 cm}
\opt{D}{30 cm}
\end{qoptions}
\qn{10} Thickness of lens is
\begin{qoptions}
\opt{A}{10 cm}
\opt{B}{15 cm}
\opt{C}{20 cm}
\opt{D}{25 cm}
\end{qoptions}
\sect{SECTION-II}
\qn{1} In given circuit reading of ammeter and voltameter are $\dfrac23$A and $\dfrac{20}{3}$V respectively. If resistance of voltameter is $R_V$ and resistance of ammeter is $R_A$ then find $R_V-R_A$ (in $\Omega$).
\qfig{p1_phy_s3_q01_1}
\qn{2} A uniform cloud of dust particle approach a planet from a great distance with velocity $v_0$ relative to the planet. Radius of planet is $R$ and escape velocity from surface of the planet is $v_e=2v_0$. Length of the cloud along direction of motion is $\ell=R$, Cross-section of cloud perpendicular to it's direction of motion is very large and it's density is $\dfrac{4}{3000}$ times density of planet. If mass of planet $m$ is very large compared to the mass of cloud then mass of the dust particle that planet will collect during passage through the cloud is $\dfrac{m}{10x}$. The value of $x$ is :
\qn{3} A plane mirror $M_1$ is placed at angle $30^\circ$ with $x$-axis and a lens $L_1$ of focal length 15 cm is placed 30 cm away from it as shown in figure. A source $O$ is placed at $[0,\,10\sqrt3\text{cm}]$. If co-ordinate of image formed after reflection from plane mirror then refraction from lens is $(x,y)$ then value of $\sqrt3\dfrac xy$ is.
\qfig{p1_phy_s3_q03_1}
\qn{4} A square non-conducting sheet of side $a$ has uniform charge density $\sigma$. A point charge $Q$ is placed at perpendicular distance $d=a$ above one of it's edges. Then component of electric force acting on sheet perpendicular to it's plane is $F$. Find $\dfrac{Q\sigma}{F\varepsilon_0}$.
\qn{5} In an infinite grid of triangular cell, resistance of each branch is $R=210\Omega$. The equivalent resistance of circuit between the terminal $A$ and $B$ is $x\Omega$. Find $x$.
\qfig{p1_phy_s3_q05_1}
\qn{6} A small block having electric dipole moment $p$ is placed on a frictionless horizontal table is attached to one end of a spring of force constant $k$ as shown. Now a point charge $-q$ is brought very slowly from a great distance towards block along line coinciding the axis of the spring. Extension in spring when block stick to charge is $\left[\dfrac{apq}{b\pi\varepsilon_0k}\right]^{\frac14}$ where $a$ and $b$ are smallest possible integer than $\dfrac{a-b}{100}$ is.
\qfig{p1_phy_s3_q06_1}
\qn{7} A ball mass $m=1$kg and radius $R=10$cm is rotating with angular velocity $\omega=20$ rad/s about horizontal axis. It dropped from a height $h=1.1$m on a rough horizontal surface ($\mu_k=0.5$) as shown in figure. If coefficient of resistitution for collision is $e=0.5$ then find horizontal distance travelled by ball between points of $1^{\text{st}}$ and $2^{\text{nd}}$ collision with ground (in m)
\qfig{p1_phy_s3_q07_1}
\qn{8} A monochromatic light ray incident on block $ABCD$ having R.I $\mu=\dfrac{2}{\sqrt3}$ at angle of incidence $i=\sin^{-1}\left[\dfrac{1}{\sqrt3}\right]$ as shown in figure. Find angle between incident and emergent ray (in degree)
\qfig{p1_phy_s3_q08_1}
\raggedcolumns
\columnbreak
\partsec{2}{CHEMISTRY}
\sect{SECTION-I(i)}
\qn{1} Which of the following is(are) correct ?
\begin{qoptions}
\opt{A}{Number of geometrical isomers of ionisation isomers of $[\mathrm{Co(NH_3)_3(H_2O)Cl_2}]\mathrm{Br}$ is 5 (except the given compound)}
\opt{B}{In $\mathrm{Fe_2(CO)_9}$, $d_{\mathrm{C-O}}$ (bridging) $<$ $d_{\mathrm{C-O}}$ (terminal)}
\opt{C}{IUPAC name of $[(\mathrm{Cl_3Sn})_2\mathrm{RhCl_2Rh(SnCl_3)_2}]^{4-}$ is Di-$\mu$-chloridotetrakis(trichlorostannito)dirhodium(I) ion}
\opt{D}{Complexes $[\mathrm{Co(NO_2)_6}]^{4-}$ and $[\mathrm{CoF_6}]^{3-}$ are inner orbital complex and outer orbital complex respectively}
\end{qoptions}
\qn{2} Which statement(s) can be possible for homonuclear diatomic neutral molecules formed by any $2^{\text{nd}}$ period element.
\begin{qoptions}
\opt{A}{Atleast one paramagnetic species has only $\sigma$ bonds}
\opt{B}{Atleast one diamagnetic species has only $\sigma$ bonds}
\opt{C}{Atleast one diamagnetic species has only $\pi$ bonds}
\opt{D}{Atleast one paramagnetic species has only $\pi$ bonds}
\end{qoptions}
\qn{3} 250 ml of 1M $\mathrm{I_2}$ is separately reacted with 0.5M $\mathrm{Cu_2S}$ solution, 0.5M $\mathrm{CuS}$ solution and 0.5 M $\mathrm{S_2O_3^{2-}}$ solution causing production of $\mathrm{Cu^{+2}}$ and $\mathrm{SO_4^{2-}}$ in the first two and $\mathrm{S_4O_6^{2-}}$ in the last case along with iodide ions. Which of the following option is/are correct assuming 100\% completion of the reaction. ($\mathrm{I}=127$)
\begin{qoptions}
\opt{A}{100 ml of $\mathrm{Cu_2S}$ solution will be consumed.}
\opt{B}{125 ml of $\mathrm{CuS}$ solution will be consumed.}
\opt{C}{500 ml of $\mathrm{S_2O_3^{2-}}$ solution will be consumed.}
\opt{D}{Equivalent weight of $\mathrm{I_2}$ in each of the reaction will be 127.}
\end{qoptions}
\qn{4} 1 mole each of benzene and toluene are mixed (ideal solution) at a given temperature. A total of 1 mole has vaporised, at given pressure. Then, which of the following must be true for this equilibrium mixture? [Given : $P^0_{\text{Benzene}}=x$; $P^0_{\text{Toluene}}=y$]. Assume $X$, $U$, $V''$, $U'$, $Z$, $V'$, $Z'$ are various points on the curve.
\begin{qoptions}
\opt{A}{Total vapour pressure $<\left[\dfrac{x+y}{2}\right]$}
\opt{B}{Total vapour pressure $=\sqrt{xy}$}
\opt{C}{In the graph,\newline \qfigi{p1_chem_s1_q04_1}\newline the vapour pressure of equilibrium mixture may correspond to point $V''$.}
\opt{D}{If in the graph given in option(C), external pressure is isothermally increased, then last trace of liquid will disappear at point $Z$.}
\end{qoptions}
\qn{5} Which of the following reaction(s) produce aldehyde as a product :
\begin{qoptions}
\opt{A}{\qfigi{p1_chem_s1_q05_A}}
\opt{B}{\qfigi{p1_chem_s1_q05_B}}
\opt{C}{$\mathrm{CH_3-CN}\ \xrightarrow[\text{(ii) }\mathrm{H_3O^+}]{\text{(i) }\mathrm{SnCl_2}/\mathrm{HCl}}$}
\opt{D}{$\mathrm{CH_3MgBr+CH_3-C\equiv N}\ \xrightarrow{\mathrm{H^+/H_2O}}$}
\end{qoptions}
\qn{6} Which reaction(s) among the following is not represent correct product formation :
\begin{qoptions}
\opt{A}{\qfigi{p1_chem_s1_q06_A}}
\opt{B}{\qfigi{p1_chem_s1_q06_B}}
\opt{C}{\qfigi{p1_chem_s1_q06_C}}
\opt{D}{\qfigi{p1_chem_s1_q06_D}}
\end{qoptions}
\sect{SECTION-I(ii)}
\paraq{7}{8}
The complex $\mathrm{HgCo.4SCN}$ has two isomers $X$ and $Y$ (having $\mathrm{Co}$ as central metal ion in both) which are HOMOLEPTIC and magnetic moment ($\mu$) of $X=Y$.
\qn{7} Choose the incorrect option for $X$ and $Y$?\newline
(I) $X$ and $Y$ are different in structure around central metal ion\newline
(II) Oxidation of Hg and Co are same in both the complexes\newline
(III) $X$ has spin only magnetic moment equals to 1.73 BM\newline
(IV) Both isomers having same hybridisation around central metal ion
\begin{qoptions}
\opt{A}{I and II}
\opt{B}{III and IV}
\opt{C}{I and III}
\opt{D}{II, III and IV}
\end{qoptions}
\qn{8} The individual sum of the linkage isomers of `X' and `Y' are.
\begin{qoptions}
\opt{A}{8}
\opt{B}{6}
\opt{C}{10}
\opt{D}{4}
\end{qoptions}
\columnbreak
\paraq{9}{10}
Alkyl halides can undergo $\mathrm{S_N1}$, $\mathrm{S_N2}$, $\mathrm{E_1}$, $\mathrm{E_2}$, $\mathrm{E_{1cb}}$ depending on various conditions.
\qfig{p1_chem_s2_para9_10_1}
\qn{9} Select the correct statement :
\begin{qoptions}
\opt{A}{Formation of $U$ follows $\mathrm{E^2}$ mechanism to give saytzeff product}
\opt{B}{Formation of $U$ follows $\mathrm{E_{1cb}}$ mechanism to give saytzeff product}
\opt{C}{Formation of $U$ follows $\mathrm{E^1}$ mechanism to give Hoffmann product}
\opt{D}{Formation of $U$ follows $\mathrm{E^2}$ mechanism to give Hoffmann product}
\end{qoptions}
\qn{10} Identify `$Y$' :
\begin{qoptions}
\opt{A}{\qfigi{p1_chem_s2_q10_A}}
\opt{B}{\qfigi{p1_chem_s2_q10_B}}
\opt{C}{\qfigi{p1_chem_s2_q10_C}}
\opt{D}{Both (B) \& (C)}
\end{qoptions}
\sect{SECTION-II}
\qn{1} Find number of process or ore given in Column-I are correctly represented with Column-II. You have to consider P for P, Q for Q, R for R etc.
\begin{center}\small\renewcommand{\arraystretch}{1.1}\setlength{\extrarowheight}{3.5pt}\setlength{\tabcolsep}{3pt}
\begin{tabular}{|c|l|c|l|}
\hline
\multicolumn{2}{|c|}{\parbox[t]{3.4cm}{\centering (I)\\ Process/Ore}} & \multicolumn{2}{c|}{\parbox[t]{3.4cm}{\centering (II)\\ Related to}}\\[3pt]
\hline
(P) & Matte smelting & (P) & Cu\\[3pt]
(Q) & Cyanide leaching & (Q) & Au\\[3pt]
(R) & Mond's process & (R) & \parbox[t]{2.3cm}{Reaction of metal $+\mathrm{H_2}$}\\[3pt]
(S) & Fused oxide electrolysis & (S) & Al\\[3pt]
(T) & Chromite ore & (T) & Fe $+$ Cr\\[3pt]
(U) & Van Arkel process & (U) & Zr\\[3pt]
(V) & Bacterial leaching & (V) & Low grade Cu ore\\[3pt]
(W) & Separation of $\mathrm{FeWO_4}$ & (W) & Extraction of Pb\\[3pt]
\hline
\end{tabular}
\end{center}
\qn{2} Consider the following statements\newline
(1) $\mathrm{CO_2}$ (g) $+$ C (s) $\to$ 2CO(g) ; it is an exothermic reaction\newline
(2) Ilmenite is an carbonate ore containing two metals\newline
(3) In froth floatation process KCN act as an activator when galena contains impurity of ZnS\newline
(4) Liquation is applicable for the purification of Sn , Pb, Zn\newline
(5) Au $+$ NaCN $+$ $\mathrm{O_2}$ $+$ $\mathrm{H_2O}$ $\to$ $\mathrm{Na[Au(CN)_4]}$ reaction occurs during leaching of gold\newline
(6) $\mathrm{Al_2O_3}$ can be obtained in solution from red bauxite by using concentrated NaOH\newline
(7) MgO was reduced by C at very high temperature followed by a rapid cooling with the help of inert gas\newline
The number of correct statement(s) in the above given statements are :
\qn{3} For the reaction , $4\mathrm A\to3\mathrm B$, the following graph is obtained experimentally, the rate of reaction (in M/sec) when $[\mathrm A]=2$M is
\qfig{p1_chem_s3_q03_1}
\qn{4} The decomposition of compound $P$, at temperature $T$ according to the equation $2P_{(g)}\to4Q_{(g)}+R_{(g)}+S_{(\ell)}$ is the first order reaction. After 30 minutes from the start of decomposition in a closed vessel, the total pressure developed is found to be 317 mm Hg and after a long period of time the total pressure observed to be 617 mm Hg. The total pressure of the vessel after 75 minute, if the volume of liquid $S$ is supposed to be negligible is: (Given : Vapour pressure of $S(\ell)$ at temperature $T=32.5$ mm Hg, $\ln(1.169)=0.16$, $\ln(1.5)=0.40$)
\qn{5} The saturated solution of a certain salt show elevation in boiling point equal to $3.12\times10^{-3}$ K in aqueous solution. If the solubility of salt is $2\times10^{-3}$ mole/kg at boiling point, then how many moles of ions are given by 1 mole of solid salt on complete dissociation.\\
$K_{b(\mathrm{H_2O})}=0.52$ K kg/mole
\qn{6} If all carbonyls isomers of molecular mass $=86$ u are separately treated with $\mathrm{CH_3MgBr}$ followed by acid hydrolysis, how many of them will give racemic mixture ?
\qn{7} How many of reagents can carry out following transformation :
\qfig{p1_chem_s3_q07_1}
(a) Cu, 300$^\circ$C\newline
(b) $\mathrm{O_3}$ / Zn, $\mathrm{H_2O}$\newline
(c) Zn--Hg / HCl\newline
(d) Jones reagent $\mathrm{CrO_3+H_2O}$\newline
(e) $\mathrm{KMnO_4}$, NaOH, $\Delta$\newline
\qfigi{p1_chem_s3_q07_2}\newline
(g) Collin's reagent\newline
(h) (i) TsCl, (ii) DMSO, $\mathrm{NaHCO_3}$ , $\Delta$\newline
(i) $\mathrm{C_5H_5HCrO_3Cl^-}$ (PCC)
\qn{8} An ideal crystal $AB$ has rock salt structure in which $A^+$ is doped with atoms of `$C$' which occupy all the tetrahedral voids without distorting the lattice. If the crystal shows defect such that the body-centered atoms is missing, then the percentage of body diagonal covered by ions/atoms in the defective crystal is (Given $\sqrt{1.5}=29/24$)
\raggedcolumns
\partsec{3}{MATHEMATICS}
\sect{SECTION-I(i)}
\qn{1} Consider $P(x)=9x^8-18x^5-20x^3+15$, then $P(x)=0$ has
\begin{qoptions}
\opt{A}{atleast one real root in $[0,3^{1/3}]$}
\opt{B}{atleast one real root in $[5^{1/5},3^{1/3}]$}
\opt{C}{atleast one real root in $[0,5^{1/5}]$}
\opt{D}{no real root in $[0,3^{1/3}]$}
\end{qoptions}
\qn{2} Let $f(x)$ be a cubic polynomial which has local maximum at $x=-1$ and $f'(x)$ has local minimum at $x=1$. If $f(-1)=10$ and $f(0)=5$ and roots of $f(x)$ are $\alpha,\beta,\gamma$, then $[\alpha]^2+[\beta]^2+[\gamma]^2$ is (where $[.]$ denotes greatest integer function)
\begin{qoptions}
\opt{A}{20}
\opt{B}{22}
\opt{C}{25}
\opt{D}{31}
\end{qoptions}
\qn{3} Let $f(x,y)=\dfrac{1}{\displaystyle\lim_{\alpha\to0}\left(\dfrac{\sqrt[x]{\sec\alpha}-\sqrt[y]{\sec\alpha}}{\alpha^2}\right)}$, then
\begin{qoptions}
\opt{A}{$f(2,3)=12$}
\opt{B}{$f(x,y)=12$ has infinitely many solutions}
\opt{C}{Slope of tangent to the curve $f(x,y)=12$ at $x=2$ is 4}
\opt{D}{$f(x,y)=12$ has no horizontal tangent}
\end{qoptions}
\qn{4} If $y^x-x^y=1$, then $\dfrac{dy}{dx}$ at $x=1$, is -
\begin{qoptions}
\opt{A}{$2(1-\ln2)$}
\opt{B}{$2(1+\ln2)$}
\opt{C}{$\ln(4e^2)$}
\opt{D}{$\ln\left(\dfrac{e^2}{4}\right)$}
\end{qoptions}
\qn{5} Identify the CORRECT statements -
\begin{qoptions}
\opt{A}{Sum of two periodic functions may or maynot be a periodic function}
\opt{B}{A differentiable function in $R$ will always be continuous in $R$}
\opt{C}{A continuous function whose domain is in a closed interval is always bounded}
\opt{D}{In $[a,b]$ for a function $f(x)$, if $f(a).f(b)<0$, then there exists atleast one root of $f(x)=0$ in $(a,b)$}
\end{qoptions}
\qn{6} If $n\in N$ and $f(x)=({}^nC_0x+1)({}^nC_1x+2)({}^nC_2x+3)({}^nC_3x+4)\ldots({}^nC_nx+n+1)$ and $f'(0)=\dfrac32(7!)$, then $n$ is equal to
\begin{qoptions}
\opt{A}{3}
\opt{B}{4}
\opt{C}{5}
\opt{D}{6}
\end{qoptions}
\sect{SECTION-I(ii)}
\paraq{7}{8}
Let $y=f(x)$ be a fourth degree algebraic polynomial function such that\newline
(i) $\displaystyle\lim_{x\to1}\dfrac{f(x)}{(x-1)^2}=2$\newline
(ii) It cuts $y$-axis a point where slope of normal is $\dfrac18$\newline
(iii) Tangent at $x=2$ is parallel to $y=8x+1$\newline
Consider $y=g(x)$ be a quadratic curve touching and cutting $y=f(x)$ on $x$-axis and $y$-axis respectively.
\qn{7} The angle of intersection between $y=f(x)$ \& $y=g(x)$ at point where they cut $y$-axis is -
\begin{qoptions}
\opt{A}{$\tan^{-1}\dfrac{2}{33}$}
\opt{B}{$\tan^{-1}\dfrac{2}{49}$}
\opt{C}{$\tan^{-1}\dfrac{14}{33}$}
\opt{D}{$\tan^{-1}\dfrac{14}{49}$}
\end{qoptions}
\qn{8} The abscissae of point on curve $y=g(x)$, where length of subtangent is 2018 can be -
\begin{qoptions}
\opt{A}{1010}
\opt{B}{2019}
\opt{C}{3028}
\opt{D}{4037}
\end{qoptions}
\paraq{9}{10}
Let $f(x)$ be a polynomial function satisfying the following conditions $\displaystyle\lim_{x\to\infty}\dfrac{f(x)}{|x|^3}=0$, $\displaystyle\lim_{x\to\infty}\left(\sqrt{f(x)}-x\right)=-1$ and $f(0)=0$ (Note : $[k]$ denotes greatest integer function less than or equal to $k$)\newline
On the basis of above information, answer the following questions :
\qn{9} The number of points of discontinuity of $g(x)=[f(x)]$ in $[0,3]$ is
\begin{qoptions}
\opt{A}{2}
\opt{B}{3}
\opt{C}{4}
\opt{D}{5}
\end{qoptions}
\qn{10} The number of points of non-derivability of $h(x)=|f(|x|)|$ is -
\begin{qoptions}
\opt{A}{2}
\opt{B}{3}
\opt{C}{4}
\opt{D}{5}
\end{qoptions}
\sect{SECTION-II}
\qn{1} The number of tangents to the curve $y=\cos(x+2y)$, $-2\pi\leqslant x\leqslant2\pi$, that are parallel to the line $x+3y=0$ is
\qn{2} Let $f(x)$ be a differentiable function on $[0,8]$ such that $f(1)=5$, $f(2)=\dfrac12$, $f(3)=6$, $f(4)=-3$, $f(5)=4$, $f(6)=\dfrac14$ and $f(7)=\dfrac{-1}{5}$. If the minimum number of roots of the equation $f'(x)-f'(x)(f(x))^2=0$ is $\lambda$, then $\dfrac\lambda3$ equals
\qn{3} Let $f(x)$ be a real valued function satisfying $2f(xy)=(f(x))^y+(f(y))^x\ \forall x,y\in R$, $f(1)=3$ and $\displaystyle\lim_{x\to0}\dfrac{f(x)+f(-x)-f(f(0))+1}{x^2}$ is equal to $(\ln k)^2$ where $k$ is natural number, then the value of $\dfrac k2$ is
\qn{4} $\displaystyle\lim_{x\to\infty}x^3\left(\sqrt{x^2+\sqrt{1+x^4}}-\sqrt2x\right)=\dfrac{\sqrt2}{\lambda}$ find value of $\lambda$
\qn{5} Let $f:(0,\infty)\to(0,\infty)$ be a function differentiable at 3 and satisfying $f(3)=3f'(3)$. If the value of $\displaystyle\lim_{x\to\infty}\left(\dfrac{f\left(3+\frac3x\right)}{f(3)}\right)^x$ is $L$ then $\dfrac{[L]}{5}$ is (where $[.]$ is greatest integer function)
\qn{6} Let $f(x)=(x-1)^{100}(x-2)^{2(99)}(x-3)^{3(98)}\ldots(x-100)^{100}$. If $P=\dfrac{f'(101)}{f(101)}$, then $\dfrac{P}{5000}$ equals
\qn{7} If a rectangular tank has a sliding panel $S$ that divides it into two adjustable tanks of width 3 ft. See figure. Water is poured into the left compartment at the rate of 5 ft$^3$/min. At the same time $S$ is moved to the right at the rate of 3ft/min. When the left comparment is 10 ft long it contains 70ft$^3$ of water, then how fast the water level is falling
\qfig{p1_math_s3_q07_1}
\qn{8} If $f(x)+f\left(1-\dfrac1x\right)=1+x\ \forall\,x\in R-\{0,1\}$, then value of $f(2)$ is
\newpage
\papersec{2}
\partsec{1}{PHYSICS}
\sect{SECTION-I(i)}
\qn{1} A conducting shell have total charge 6C is placed inside a uniform EF of 80 N/C as shown. Net force on 4C charge particle placed inside shell is $-80\,\mathrm N\,\hat i$ and conductor is held fixed then select the incorrect options :-
\qfig{p2_phy_s1_q01_1}
\begin{qoptions}
\opt{A}{Force on 4C due to charge induced on outer surface is $80\,\mathrm N\,\hat i$.}
\opt{B}{Force on 4C due to charge induced on outer surface is $320\,\mathrm N\,\hat i$.}
\opt{C}{Force on shell due to 4C is $400\,\mathrm N\,\hat i$.}
\opt{D}{Force on charge particle due to charge induced on inner surface of shell is $+80\,\mathrm N\,\hat i$}
\end{qoptions}
\qn{2} Consider electric circuit shown in figure. All batteries are of emf 2V and all resistors are of $1\Omega$. Choose correct options :-
\qfig{p2_phy_s1_q02_1}
\begin{qoptions}
\opt{A}{Current in resistor between $C$ \& $G$ is $\dfrac87$A.}
\opt{B}{$V_A-V_E$ is equal to $\dfrac{49}{7}$V.}
\opt{C}{$V_B-V_F$ is equal to $\dfrac{20}{7}$V.}
\opt{D}{Battery between $C$ \& $D$ supplies power of $\dfrac47$W.}
\end{qoptions}
\qn{3} For an equilateral glass prism graph of deviation $\delta$ versus angle of incidence $i$ is shown in figure. Choose incorrect option(s) :-
\qfig{p2_phy_s1_q03_1}
\begin{qoptions}
\opt{A}{Maximum possible deviation produced by prism is $60^\circ+\sin^{-1}\left(\sqrt{\dfrac32}-\dfrac12\right)$}
\opt{B}{Maximum possible deviation produced by prism is $30^\circ+\sin^{-1}\left(\sqrt{\dfrac32}+\dfrac12\right)$.}
\opt{C}{Refraction index of prism is $\sqrt{\dfrac32}$.}
\opt{D}{For normal incidence on prism total deviation suffered by ray is $60^\circ$.}
\end{qoptions}
\qn{4} Charge density on right half of an insulating spherical shell is positive and charge density on left half is negative but same magnitude as on left half. Choose correct options :-
\qfig{p2_phy_s1_q04_1}
\begin{qoptions}
\opt{A}{Electric field on $x$-axis between $A$ \& $B$ is in same direction at all points.}
\opt{B}{Electric field between $A$ \& $O$ and $O$ \& $B$ are along opposite directions.}
\opt{C}{$y$-$z$ plane inside sphere is equipotential.}
\opt{D}{Electric field on all points on $x$-axis is in same direction outside the sphere.}
\end{qoptions}
\qn{5} An image of a virtual object is formed by a spherical mirror and height of image $I$ is half of height of object $O$ from principal axis. Select correct options :-
\qfig{p2_phy_s1_q05_1}
\begin{qoptions}
\opt{A}{Mirror is on the right side of $I$ and it is convex.}
\opt{B}{Focal length of mirror is 20 cm.}
\opt{C}{Focal length of mirror is $\dfrac{80}{3}$ cm.}
\opt{D}{Mirror is concave and on right side of object.}
\end{qoptions}
\qn{6} An uniform insulating ring and radius $R$ has total charge $q$ distributed uniformly along its circumference except on length $d$ ($d\ll R$). It is placed on smooth surface in uniform electric field parallel to length $d$ which is between $A$ \& $B$. Initially ring is released from rest. Select correct options :-
\qfig{p2_phy_s1_q06_1}
\begin{qoptions}
\opt{A}{Initial torque about COM is $\dfrac{qEd}{2\pi}$.}
\opt{B}{Initial torque about COM is $\dfrac{qEd}{\pi}$}
\opt{C}{Maximum angular speed in subsequent motion will be $\dfrac1R\sqrt{\dfrac{qdE}{2\pi m}}$.}
\opt{D}{Maximum angular speed in subsequent motion will be $\dfrac1R\sqrt{\dfrac{qdE}{\pi m}}$.}
\end{qoptions}
\sect{SECTION-I(ii)}
\qn{7} Between two concentric conducting neutral shells, a charge $Q=1$C is fixed at distance of 2m from centre $O$. Potential at centre is $V_0$, at inner sphere is $V_1$ and outer sphere is $V_2$ and total charge on outer surfaces of inner \& outer spheres are $Q_1$ and $Q_2$ respectively. After connecting both with conducting wire these charges becomes $Q_1'$ and $Q_2'$. Match List--I with List--II.
\qfig{p2_phy_s2_q07_1}
\begin{center}\small\renewcommand{\arraystretch}{1.1}\setlength{\extrarowheight}{3.5pt}\setlength{\tabcolsep}{4pt}
\begin{tabular}{|c|c|c|c|}
\hline
\multicolumn{2}{|c|}{List-I} & \multicolumn{2}{c|}{List-II}\\[3pt]
\hline
(P) & $\dfrac{V_1}{V_2}$ & (1) & 3\\[6pt]
(Q) & $\left|\dfrac{Q_1}{Q_2}\right|$ & (2) & 1\\[6pt]
(R) & $\left|\dfrac{Q_1'}{Q_2'}\right|$ & (3) & 2\\[6pt]
(S) & $\dfrac{V_0}{V_2}$ & (4) & $\dfrac13$\\[6pt]
 & & (5) & 0\\[3pt]
\hline
\end{tabular}
\end{center}
\begin{qoptions}
\opt{A}{$P\to3;Q\to5;R\to4;S\to3$}
\opt{B}{$P\to2;Q\to2;R\to5;S\to2$}
\opt{C}{$P\to5;Q\to5;R\to5;S\to2$}
\opt{D}{$P\to1;Q\to5;R\to5;S\to1$}
\end{qoptions}
\qn{8} A ball of radius $r=(\sqrt2-1)R$ falls on both balls symmetrically as shown with speed 10 m/s. If $A$ \& $B$ moves with 3 m/s just after collision then match the List--I with list--II. All quantities in List--II are in SI units.
\qfig{p2_phy_s2_q08_1}
\begin{center}\small\renewcommand{\arraystretch}{1.1}\setlength{\extrarowheight}{3.5pt}\setlength{\tabcolsep}{4pt}
\begin{tabular}{|c|l|c|c|}
\hline
 & \multicolumn{1}{c|}{List-I} & & List-II\\[3pt]
\hline
(P) & \parbox[t]{4.2cm}{Magnitude of impulse of normal between $A$ \& $C$.} & (1) & 6\\[3pt]
(Q) & \parbox[t]{4.2cm}{Magnitude of impulse of tension in any string.} & (2) & 0.2\\[3pt]
(R) & \parbox[t]{4.2cm}{Velocity of $C$ just after collision.} & (3) & $6\sqrt2$\\[3pt]
(S) & \parbox[t]{4.2cm}{Magnitude of net impulse on $C$ during collision.} & (4) & 12\\[3pt]
 & & (5) & 2\\[3pt]
\hline
\end{tabular}
\end{center}
\begin{qoptions}
\opt{A}{$P\to1;Q\to3;R\to1;S\to2$}
\opt{B}{$P\to3;Q\to1;R\to5;S\to4$}
\opt{C}{$P\to1;Q\to1;R\to2;S\to4$}
\opt{D}{$P\to1;Q\to1;R\to2;S\to2$}
\end{qoptions}
\qn{9} A long thin carpet of length 4m is being laid on the floor. One end of the carpet is bent and then pulled backward with constant velocity $v=2$m/s above the part of carpet which is still at rest on the floor. Carpet has mass density of 0.5 kg/m. All quantities in List--II in SI units.
\qfig{p2_phy_s2_q09_1}
\begin{center}\small\renewcommand{\arraystretch}{1.1}\setlength{\extrarowheight}{3.5pt}\setlength{\tabcolsep}{4pt}
\begin{tabular}{|c|l|c|c|}
\hline
 & \multicolumn{1}{c|}{List-I} & & List-II\\[3pt]
\hline
(P) & \parbox[t]{4.2cm}{Velocity of com of part of carpet still at rest.} & (1) & 2\\[3pt]
(Q) & \parbox[t]{4.2cm}{Velocity of com of moving part of carpet.} & (2) & 1.5\\[3pt]
(R) & \parbox[t]{4.2cm}{Mass being added in moving part per sec.} & (3) & 1\\[3pt]
(S) & \parbox[t]{4.2cm}{Force needed to pull the moving part.} & (4) & 0.5\\[3pt]
 & & (5) & 0\\[3pt]
\hline
\end{tabular}
\end{center}
\begin{qoptions}
\opt{A}{$P\to5;Q\to2;R\to3;S\to1$}
\opt{B}{$P\to4;Q\to2;R\to4;S\to1$}
\opt{C}{$P\to5;Q\to2;R\to4;S\to3$}
\opt{D}{$P\to3;Q\to3;R\to3;S\to1$}
\end{qoptions}
\qn{10} There is a point source of light at a depth 8m below the surface of water. An observer in air sees the image of the source. Line of sight of observer makes an angle $53^\circ$ with the normal to surface. All quantities in List--II are in SI unit. (use : $\sqrt7=2.65$, $\mu$ of water $=4/3$) ?
\begin{center}\small\renewcommand{\arraystretch}{1.1}\setlength{\extrarowheight}{3.5pt}\setlength{\tabcolsep}{4pt}
\begin{tabular}{|c|l|c|c|}
\hline
 & \multicolumn{1}{c|}{List-I} & & List-II\\[3pt]
\hline
(P) & \parbox[t]{4.2cm}{Apparent depth of image from surface.} & (1) & 9.05\\[3pt]
(Q) & \parbox[t]{4.2cm}{Horizontal distance between object and its image.} & (2) & 12.02\\[3pt]
(R) & \parbox[t]{4.2cm}{Radius of circle of illuminance.} & (3) & 3.38\\[3pt]
(S) & \parbox[t]{4.2cm}{Percentage of power coming out of the water assuming no absorption of light by water and 100\% refraction within circle of illuminance.} & (4) & 16.34\\[3pt]
 & & (5) & 2.53\\[3pt]
\hline
\end{tabular}
\end{center}
\begin{qoptions}
\opt{A}{$P\to3;Q\to1;R\to5;S\to2$}
\opt{B}{$P\to2;Q\to1;R\to4;S\to3$}
\opt{C}{$P\to4;Q\to5;R\to1;S\to4$}
\opt{D}{$P\to5;Q\to3;R\to1;S\to4$}
\end{qoptions}
\sect{SECTION-II}
\qn{1} A light ray incident at graving incidence on a long rectangular slab having thickness $d$ at origin. Determine $x$ coordinate of point $A$ where ray meets the upper surface of slab air interface. $\mu$ of slab varies as $\mu=\sqrt{1+e^{x/d}}$; $d=4$m and take $\ln2=0.7$
\qfig{p2_phy_s3_q01_1}
\qn{2} Consider a uniformly charged very long solid non-conducting cylinder of radius $R$ and having volume charge density $\rho=\epsilon_0$. Find electric flux (in SI unit) through a gaussian spherical surface of radius $2R$ and having one of its diameter along the axis of cylinder and centre on axis of cylinder. Take : $R=\left(\dfrac1\pi\right)^{1/3}$
\qn{3} Bottom of a container is made of an equi-convex lens having focal length 60 cm in air and $\mu=1.5$. At what height (cm) above water surface an object should be placed so that final image of object coincides with object. Lower surface of lens is silvered. ($\mu_{\text{water}}=4/3$)
\qfig{p2_phy_s3_q03_1}
\qn{4} A rocket is launched from point $A$, which is at distance $4R$ from centre of earth. If it passes grazing the earth surface then find value of $\dfrac{v}{100}$ where $v$ is velocity of projection. $\left[\text{use }\dfrac{GM}{R}=6.4\times10^7\,\mathrm{m^2/s^2}\right]$
\qfig{p2_phy_s3_q04_1}
\qn{5} A hollow hemispherical bowl of mass $m$ is placed on a smooth surface as shown. A particle of same mass strikes it and sticks with the bowl at $A$. Find velocity of particle just after collision in vertical direction ? ($u=5$ m/s)
\qfig{p2_phy_s3_q05_1}
\qn{6} A beam of fast moving electrons having cross-sectional area $A=1$ cm$^2$ falls normally on a flat surface. The electrons are absorbed by the surface and the average pressure exerted by the electrons on this surface is found to be $P=9.1$ Pa. If the electrons are moving with a speed $v=24\times10^7$ m/s, then find the effective current (in A) through any cross-section of the electron beam. (mass of electron $=9.1\times10^{-31}$ kg)
\qn{7} A meteorite is approaching a planet in a straight line passing through the centre of planet. It collides with a satellite orbiting the planet in orbit of radius $r$ having mass ten times of mass of meteorite and sticks with it. As a result of collision satellite goes to new orbit with minimum distance $\dfrac r2$ from the centre of planet. If speed of meteorite just before collision is $u=\sqrt{\dfrac{nGM}{R}}$. Find value of $n$.
\qn{8} Figure shows a binary star system revolving about their COM. The masses of star $A$ \& $B$ are $15\times10^{30}$ kg and $30\times10^{30}$ kg respectively. Find the ratio of area swept by star $A$ to area swept by star $B$ in a common time interval.
\qfig{p2_phy_s3_q08_1}
\raggedcolumns
\columnbreak
\partsec{2}{CHEMISTRY}
\sect{SECTION-I(i)}
\qn{1} Element ``$X$'' forms a compound $\mathrm{XCl_2.4NH_3}$ that does not show any type of isomerism. Element ``$Y$'' forms $\mathrm{YCl_3.4NH_3}$ type compound which has two isomers ``$P$'' \& ``$Q$''. Also, both elements ``$X$'' \& ``$Y$'' belong to VIII group of Mendeleev's periodic table.\newline
``$Y$'' is the only metal present in vitamin $\mathrm{B_{12}}$ and compounds of ``$X$'' inhibit growth of tumours.\newline
Which is/are correct about $X$ \& $Y$ or their compounds?
\begin{qoptions}
\opt{A}{If one $\mathrm{NH_3}$ is removed from the formula of $\mathrm{YCl_3.4NH_3}$, then one of the isomer of newly formed complex compound becomes optically active}
\opt{B}{Both ``$P$'' \& ``$Q$'' are $1:1$ electrolytes}
\opt{C}{If two $\mathrm{NH_3}$ in $\mathrm{XCl_2.4NH_3}$ get replaced by Cl then the formula of the new compound will show stereoisomerism}
\opt{D}{$X$ and $X^{2+}$ have same number of unpaired electrons}
\end{qoptions}
\qn{2} Which of the following statements(s) about extraction of Fe is/are correct?
\begin{qoptions}
\opt{A}{Cast iron is made by mixing pig iron with scrap iron and coke using air blast at 300K}
\opt{B}{By oxidizing impurities of cast iron in a furnace lined with Haematite, wrought iron can be produced}
\opt{C}{Layer of Molten slag prevents further oxidation state of Molten metal layer in Blast furnace}
\opt{D}{Pig iron has more carbon content than that of cast iron}
\end{qoptions}
\qn{3} Identify which of the following solution(s) will show negative deviation from Raoult's law -
\begin{qoptions}
\opt{A}{Carbon-disulphide $+$ acetic acid}
\opt{B}{Water $+$ Formic Acid}
\opt{C}{Propan-2-ol $+$ 2-methyl propanol}
\opt{D}{Chloroform $+$ Acetone}
\end{qoptions}
\qn{4} Pure benzene has a vapour pressure 900 torr and pure toluene has a vapour pressure 360 torr. If a solution is prepared by mixing 39gm benzene and 46 gm toluene then
\begin{qoptions}
\opt{A}{Vapour pressure of its distillate will be be 745.714 torr}
\opt{B}{Vapour pressure of initial solution will be 630 torr}
\opt{C}{Mole \% of benzene in vapour phase of solution will be $\dfrac{500}{7}\%$}
\opt{D}{Mole \% of toluene in its distillate will be 50\%.}
\end{qoptions}
\qn{5} Which of the following reaction will go faster if the concentration of the nucleophile / base is increased?
\begin{qoptions}
\opt{A}{\qfigi{p2_chem_s1_q05_A}}
\opt{B}{\qfigi{p2_chem_s1_q05_B}}
\opt{C}{$\mathrm{(CH_3)_3CCl}\ \xrightarrow[\mathrm{CH_3COOH}]{\mathrm{H_2O}}$}
\opt{D}{\qfigi{p2_chem_s1_q05_D}}
\end{qoptions}
\qn{6} Choose the correct statement(s) regarding following reaction :
\qfig{p2_chem_s1_q06_1}
\begin{qoptions}
\opt{A}{$(Q)$, will give only 1 product with NBS}
\opt{B}{$(R)$, is optically inactive due to internal compensation}
\opt{C}{Product $(Q)$ is trans-but-2-ene}
\opt{D}{$(P)$, can exist in four stereoisomeric forms.}
\end{qoptions}
\sect{SECTION-I(ii)}
\qn{7} Match the List.
\begin{center}\footnotesize\renewcommand{\arraystretch}{1.1}\setlength{\extrarowheight}{3.5pt}\setlength{\tabcolsep}{2pt}
\begin{tabular}{|c|l|c|l|}
\hline
\multicolumn{2}{|c|}{\parbox[t]{2.8cm}{\centering List-I\\ Complex Compound/Ion}} & \multicolumn{2}{c|}{\parbox[t]{2.8cm}{\centering List-II\\ Characteristics}}\\[3pt]
\hline
(P) & $[\mathrm{Zn(gly)_2}]$ & (I) & \parbox[t]{2.6cm}{Exhibits Geometrical Isomerism only}\\[3pt]
(Q) & $[\mathrm{Pd(NH_2CH(CH_3)COO^-)_2}]$ & (II) & \parbox[t]{2.6cm}{Exhibits Optical isomerism only}\\[3pt]
(R) & $[\mathrm{Co(C_2S_2O_2)_2(NH_3)_2}]^-$ & (III) & \parbox[t]{2.6cm}{Both geometrical \& optical isomerism are exhibited}\\[3pt]
(S) & $[\mathrm{Cu(NH_3)_4}][\mathrm{CuCl_4}]$ & (IV) & \parbox[t]{2.6cm}{Structural isomerism is exhibited}\\[3pt]
 & & (V) & \parbox[t]{2.6cm}{Square planar around any one central atom}\\[3pt]
 & & (VI) & \parbox[t]{2.6cm}{Tetrahedral around any one central atom}\\[3pt]
\hline
\end{tabular}
\end{center}
Which of the following option is correct match?
\begin{qoptions}
\opt{A}{$P\to VI;Q\to III;R\to III;S\to VI$}
\opt{B}{$P\to IV;Q\to V;R\to IV;S\to V$}
\opt{C}{$P\to V;Q\to II;R\to VI;S\to I$}
\opt{D}{$P\to II;Q\to I;R\to I;S\to III$}
\end{qoptions}
\qn{8} Match the List.
\begin{center}\footnotesize\renewcommand{\arraystretch}{1.1}\setlength{\extrarowheight}{3.5pt}\setlength{\tabcolsep}{2pt}
\begin{tabular}{|c|l|c|l|}
\hline
\multicolumn{2}{|c|}{List-I} & \multicolumn{2}{c|}{List-II}\\[3pt]
\hline
(P) & NaCl structure & (1) & \parbox[t]{3.1cm}{Cation forming fcc and anion at octahedral void}\\[3pt]
(Q) & Zinc Blende structure & (2) & \parbox[t]{3.1cm}{Anion forming fcc and at all octahedral void while cation goes to all tetrahedral void}\\[3pt]
(R) & Ceasium Halide structure & (3) & \parbox[t]{3.1cm}{Anion forming fcc and cation at alternate tetrahedral void}\\[3pt]
(S) & Flourite structure & (4) & \parbox[t]{3.1cm}{Cation forming fcc and anion at alternate tetrahedral void.}\\[3pt]
 & & (5) & \parbox[t]{3.1cm}{Anion forming simple cube \& cation at cubical void.}\\[3pt]
 & & (6) & \parbox[t]{3.1cm}{Cation forming fcc \& anion at all tetrahedral void.}\\[3pt]
\hline
\end{tabular}
\end{center}
Which of the following option is correct match?
\begin{qoptions}
\opt{A}{$P\to1;Q\to3,4;R\to2,5;S\to6$}
\opt{B}{$P\to1;Q\to3,4;R\to2;S\to6$}
\opt{C}{$P\to3,4;Q\to1;R\to2;S\to6$}
\opt{D}{$P\to3;Q\to4;R\to2,5;S\to1$}
\end{qoptions}
\qn{9} For the reaction $2\mathrm A(g)\to\mathrm B(g)$, starting with 4 moles of $A$ \& half life 40 sec in a 5 litre container at 400 K.
\begin{center}\small\renewcommand{\arraystretch}{1.1}\setlength{\extrarowheight}{3.5pt}\setlength{\tabcolsep}{3pt}
\begin{tabular}{|c|l|c|c|}
\hline
 & \multicolumn{1}{c|}{List-I} & & List-II\\[3pt]
\hline
(P) & \parbox[t]{4cm}{If reaction follow first order then initial rate of reaction} & (1) & $3.46\times10^{-4}$ M/s\\[3pt]
(Q) & \parbox[t]{4cm}{If reaction follow zero order then rate of reaction at time : $t=40$ sec.} & (2) & $6.93\times10^{-3}$ M/s\\[3pt]
(R) & \parbox[t]{4cm}{If reaction follow first order kinetics then rate of reaction at 95\% completion.} & (3) & $5\times10^{-3}$ M/s\\[3pt]
(S) & \parbox[t]{4cm}{If reaction follow second order kinetics then rate of reaction at 99\% completion} & (4) & $10^{-6}$ M/s\\[3pt]
 & & (5) & $6.25\times10^{-5}$ M/s\\[3pt]
\hline
\end{tabular}
\end{center}
\begin{qoptions}
\opt{A}{$P\to2;Q\to3;R\to1;S\to4$}
\opt{B}{$P\to1;Q\to2;R\to3;S\to4$}
\opt{C}{$P\to5;Q\to1;R\to2;S\to4$}
\opt{D}{$P\to2;Q\to1;R\to4;S\to3$}
\end{qoptions}
\qn{10} Match the list-
\begin{center}\small\renewcommand{\arraystretch}{1.1}\setlength{\extrarowheight}{3.5pt}\setlength{\tabcolsep}{3pt}
\begin{tabular}{|c|c|c|l|}
\hline
\multicolumn{2}{|c|}{\parbox[t]{3.6cm}{\centering List-I\\ (Pair of compound)}} & \multicolumn{2}{c|}{\parbox[t]{2.4cm}{\centering List-II\\ (Differentiated by)}}\\[3pt]
\hline
(P) & \qfigi{p2_chem_s2_q10_1} & (1) & Fehling test\\[3pt]
(Q) & \qfigi{p2_chem_s2_q10_2} & (2) & Tollen's test\\[3pt]
(R) & \qfigi{p2_chem_s2_q10_3} & (3) & Na-metal\\[3pt]
(S) & \qfigi{p2_chem_s2_q10_4} & (4) & $\mathrm{NaHCO_3}$\\[3pt]
 & & (5) & Cold $\mathrm{KMnO_4}$\\[3pt]
\hline
\end{tabular}
\end{center}
\begin{qoptions}
\opt{A}{$P\to1,2,;Q\to2;R\to1,2,4;S\to1,2,5$}
\opt{B}{$P\to1,2;Q\to1,2;R\to1,2;S\to1,2$}
\opt{C}{$P\to1;Q\to1,2;R\to1,2,4;S\to1,2,5$}
\opt{D}{$P\to1,2;Q\to2;R\to1,4;S\to5$}
\end{qoptions}
\sect{SECTION-II}
\qn{1} Number of reactions in which atleast one gaseous product is formed.\newline
(i) $\mathrm{Al_2O_3+SiO_2+N_2+C}\xrightarrow{1800^\circ\mathrm{C}-2000^\circ\mathrm{C}}$\newline
(ii) $\mathrm{4Au+8NaCN+O_2+2H_2O}\rightleftharpoons$\newline
(iii) $\mathrm{SnSiO_3+CaO+C}\to$\newline
(iv) $\mathrm{2Cu_2S+3O_2}\to$\newline
(v) $\mathrm{XeF_2+H_2O}\to$\newline
(vi) $\mathrm{ZnO+C}$ (coke) $\xrightarrow{1673\mathrm K}$
\qn{2} During extraction of Magnesium by Dow's process, magnesium hydroxide is treated with HCl followed by concentration of solution and crystallisation to give crystals of $\mathrm{MgCl_2.xH_2O}$. The value of ``$x$'' is \underline{\hspace{1.5cm}}.
\qn{3} Number of specie(s) that use atleast one double dumbell shape orbital for hybridization of central atom.\newline
(i) $[\mathrm{PtCl_4}]^{2-}$\quad (ii) $\mathrm{Ni(CO)_4}$\newline
(iii) $\mathrm{PCl_5}$\quad (iv) $[\mathrm{CoF_6}]^{3-}$\newline
(v) $\mathrm{K_4[Ni(CN)_4]}$\quad (vi) $\mathrm{XeF_6}$\newline
(vii) $\mathrm{Fe(CO)_5}$\quad (viii) $\mathrm{SO_2Cl_2}$
\qn{4} A solid cube with edge length of 2.532 cm of an ionic compound which has NaCl type lattice is added to 1kg of water. The boiling point of this solution is found to be $100.52^\circ$C ( Assume 100\% ionisation of ionic compound). If radius of anion of ionic solid is 2\,\AA{}, then calculate radius of cation in \AA{}.\newline
Given : $k_b$ of water $=0.52$ K kg mol$^{-1}$,\\[2pt]
$N_A=6\times10^{23}$, $\sqrt[3]{216}=6$, $(2.532)^3\times\dfrac43=21.6$
\qn{5} In $\mathrm{CaF_2}$-type lattice, if\newline
$x=$ nearest distance between $\mathrm{Ca^{2+}}$ \& $\mathrm{F^-}$ ions\newline
$y=$ nearest distance between two $\mathrm{F^-}$ ions\newline
$z=$ nearest distance between two $\mathrm{Ca^{2+}}$ ions.\newline
then find value of the expression.
\[\left[\sqrt6\left(\dfrac{xz}{y^2}\right)\right]\]
\qn{6} Which of the following can give silver mirror with Tollen's reagent :
\qfig{p2_chem_s3_q06_1}
\qn{7} How many moles of $\mathrm{CO_2}$ are released in following reaction :
\qfig{p2_chem_s3_q07_1}
\qn{8} How many of the following are suitable reagents for following conversion :
\qfig{p2_chem_s3_q08_1}
(a) (i) $\mathrm{Ni/H_2}$ (ii) TsCl (iii) $\mathrm{NaBH_4/H_2O}$\newline
(b) $\mathrm{ZnHg/HCl}$\newline
(c) (i) $\mathrm{N_2H_4}$ / (ii) $\mathrm{\overset{\ominus}{O}H}$\newline
(d) $\mathrm{N_2H_4/H_2O_2}$\newline
(e) Na / Liq. $\mathrm{NH_3}$\newline
(f) $\mathrm{CH_3MgBr}$ followed by LAH\newline
(g) (i) DIBAL-H ($-78^\circ$C) , (ii) $\mathrm{H_2O}$
\raggedcolumns
\columnbreak
\partsec{3}{MATHEMATICS}
\sect{SECTION-I(i)}
\qn{1} If $f(x)=\begin{cases}\cos^{-1}(\cos x)\,; & |x|<\pi\\[2pt] \pi\log_\pi|x|\,; & |x|\geqslant\pi\end{cases}$, then correct statement(s) is/are -
\begin{qoptions}
\opt{A}{$f(x)$ is not differentiable at three points}
\opt{B}{$f'(\pi^+)+f'(-\pi^-)=0$}
\opt{C}{$f'(\pi^+)=f'(-\pi^-)=0$}
\opt{D}{$f'(0^+)+f'(0^-)=0$}
\end{qoptions}
\qn{2} Let the function $p,q$ and $r$ defined by $p(x)=x^{\ln x}$, $q(x)=(\ln x)^x$, $r(x)=\log_{e^x}x$, then the value of $\dfrac{p'(e^{-1}).r'(e^2)}{q'(e)}$ is (where $y'$ denotes the derivative of $y$)
\begin{qoptions}
\opt{A}{$\dfrac1e$}
\opt{B}{$\dfrac1{e^2}$}
\opt{C}{$\dfrac2e$}
\opt{D}{$\dfrac2{e^2}$}
\end{qoptions}
\qn{3} Let $f(x)$ be a cubic polynomial which is having local maximum at $(1,2)$ and $f'(x)$ has local extremum at $x=0$. If $f(0)=1$ then
\begin{qoptions}
\opt{A}{Number of real roots of $f(x)-f(-x)=6x-10$ is 1}
\opt{B}{If $g(x)+f(x)=0$ then complete interval in which $g(x)$ is increasing is $x\in(-\infty,-1]\cup[1,\infty)$}
\opt{C}{The range of values of $K$ for which $f(x)=K$ has three distinct solutions is $K\in(0,2)$}
\opt{D}{Range of the function $P(\theta)=f(\sin\theta)$ where $\theta\in R$ is $[0,2]$}
\end{qoptions}
\qn{4} For the curve $\sin x+\sin y=1$, $\displaystyle\lim_{x\to0}\left(x^\alpha\right)\dfrac{d^2y}{dx^2}$ exists then $\alpha$ can be
\begin{qoptions}
\opt{A}{$\dfrac32$}
\opt{B}{2}
\opt{C}{$\dfrac52$}
\opt{D}{1}
\end{qoptions}
\qn{5} If $f(x)=\displaystyle\lim_{n\to\infty}\dfrac{2x^{2n}+x+2}{x^{2n}-x^2+3}$, then which of the following statement(s) holds good ?
\begin{qoptions}
\opt{A}{$f(x)$ is continuous at $x=1$}
\opt{B}{$f(x)$ is discontinuous at $x=1$}
\opt{C}{$\displaystyle\lim_{x\to2}f(x)$ exists finitely}
\opt{D}{$\displaystyle\lim_{x\to2}f(x)$ does not exist}
\end{qoptions}
\qn{6} Let $a\in R$ and $f:R\to R$ be given by $f(x)=3x^5-50x^3+135x-a$. The equation $f(x)=0$ will have :
\begin{qoptions}
\opt{A}{one real root for $|a|>216$}
\opt{B}{three real roots for $88<|a|<216$}
\opt{C}{five real roots for $|a|<88$}
\opt{D}{five real roots for $a\in R$}
\end{qoptions}
\sect{SECTION-I(ii)}
\qn{7} Let $\phi(x)=3f\left(\dfrac{x^2}{3}\right)+f\left(3-x^2\right)\ \forall\,x\in(-3,4)$ where $f''(x)>0\ \forall\,x\in(-3,4)$,\\
$g(x)=(a^2-3a+2)\left(\cos^2\dfrac x4-\sin^2\dfrac x4\right)+(a-1)x+\sin1$\\
Match List-I with List-II and select the correct answer using the code given below the list.
\begin{center}\footnotesize\renewcommand{\arraystretch}{1.1}\setlength{\extrarowheight}{3.5pt}\setlength{\tabcolsep}{2pt}
\begin{tabular}{|c|l|c|l|}
\hline
\multicolumn{2}{|c|}{List-I} & \multicolumn{2}{c|}{List-II}\\[3pt]
\hline
(P) & $\phi(x)$ is increasing in & (1) & $(0,1)\cup(1,4)$\\[4pt]
(Q) & $\phi(x)$ is decreasing in & (2) & $\left[-\dfrac32,0\right]\cup\left[\dfrac32,4\right)$\\[6pt]
(R) & \parbox[t]{3.1cm}{The set of values of $a$ for which $g(x)$ does not possess critical point} & (3) & $\left(-3,-\dfrac32\right]\cup\left[0,\dfrac32\right]$\\[10pt]
(S) & \parbox[t]{3.1cm}{The set of values of $a$ for which $g(x)$ possess critical point} & (4) & $(-\infty,0]\cup\{1\}\cup[4,\infty)$\\[3pt]
\hline
\end{tabular}
\end{center}
\begin{qoptions}
\opt{A}{$P\to2;Q\to3;R\to1;S\to4$}
\opt{B}{$P\to2;Q\to3;R\to4;S\to1$}
\opt{C}{$P\to3;Q\to2;R\to1;S\to4$}
\opt{D}{$P\to3;Q\to2;R\to4;S\to1$}
\end{qoptions}
\qn{8} Match List-I with List-II and select the correct answer using the code given below the list.
\begin{center}\footnotesize\renewcommand{\arraystretch}{1.1}\setlength{\extrarowheight}{3.5pt}\setlength{\tabcolsep}{2pt}
\begin{tabular}{|c|l|c|c|}
\hline
 & \multicolumn{1}{c|}{List-I} & & List-II\\[3pt]
\hline
(P) & \parbox[t]{4.4cm}{If $x=\alpha$ satisfies $\sin^{-1}x+\cos^{-1}x^2+\dfrac\pi2=0$, then the value of $\dfrac{\sec^{-1}\alpha-\tan^{-1}\alpha}{\cot^{-1}\alpha-\mathrm{cosec}^{-1}\alpha}$ is} & (1) & 0\\[14pt]
(Q) & \parbox[t]{4.4cm}{If $a=\displaystyle\lim_{x\to\infty}\left(2.3^x+3.4^x+4.5^x\right)^{\frac1x}$ and $b=\displaystyle\lim_{x\to(-\infty)}\left(2.3^x+3.4^x+4.5^x\right)^{\frac1x}$ then $\left(\dfrac{a+b}{a-b}\right)$ is equal to} & (2) & 4\\[22pt]
(R) & \parbox[t]{4.4cm}{Number of points of discontinuity of $f(x)=x\sin(\pi(2x-[x]))$ in $[0,2]$ will be (Note $[k]$ denotes greatest integer function less than or equal to $k$)} & (3) & 1\\[22pt]
(S) & \parbox[t]{4.4cm}{$\displaystyle\lim_{x\to0}\left(\dfrac{3e^{x/3}-x-3}{x^2}\right)$ is equal to} & (4) & $\dfrac13$\\[6pt]
 & & (5) & $\dfrac12$\\[6pt]
 & & (6) & $\dfrac16$\\[3pt]
\hline
\end{tabular}
\end{center}
\begin{qoptions}
\opt{A}{$P\to1;Q\to2;R\to2;S\to4$}
\opt{B}{$P\to3;Q\to1;R\to2;S\to6$}
\opt{C}{$P\to3;Q\to2;R\to1;S\to6$}
\opt{D}{$P\to1;Q\to3;R\to2;S\to4$}
\end{qoptions}
\qn{9} Match List-I with List-II and select the correct answer using the code given below the list.
\begin{center}\footnotesize\renewcommand{\arraystretch}{1.1}\setlength{\extrarowheight}{3.5pt}\setlength{\tabcolsep}{2pt}
\begin{tabular}{|c|l|c|c|}
\hline
 & \multicolumn{1}{c|}{List-I} & & List-II\\[3pt]
\hline
(P) & \parbox[t]{4.4cm}{If $g(x)=f(-x+xf(x))$ and $f(0)=1$, $f'(0)=2$, then value of $g'(0)$ is equal to} & (1) & $\dfrac14$\\[10pt]
(Q) & \parbox[t]{4.4cm}{If $e^y=e^xx^{-y}$, then $\left(\dfrac{dy}{dx}\right)_{x=e}$ is equal to} & (2) & 6\\[8pt]
(R) & \parbox[t]{4.4cm}{Number of points of discontinuity of $y=[2x]$ in $[-1,2]$ is equal to (where $[.]$ denote greatest integer function)} & (3) & 0\\[10pt]
(S) & \parbox[t]{4.4cm}{If $I=\displaystyle\lim_{x\to0}(1+\tan x)^{\frac1{\sin x}}$ then $[I]$ is equal to (where $[.]$ denote greatest integer function)} & (4) & 2\\[10pt]
 & & (5) & 3\\[3pt]
\hline
\end{tabular}
\end{center}
\begin{qoptions}
\opt{A}{$P\to3;Q\to1;R\to2;S\to4$}
\opt{B}{$P\to1;Q\to2;R\to3;S\to5$}
\opt{C}{$P\to2;Q\to1;R\to4;S\to5$}
\opt{D}{$P\to2;Q\to1;R\to3;S\to5$}
\end{qoptions}
\qn{10} If a curve $f(x,y,a,b)=0$ is given by $\sin^{-1}x+\cos^{-1}ay+\cos^{-1}(bxy)=\dfrac\pi2$\\
Match List-I with List-II and select the correct answer using the code given below the list.
\begin{center}\footnotesize\renewcommand{\arraystretch}{1.1}\setlength{\extrarowheight}{3.5pt}\setlength{\tabcolsep}{2pt}
\begin{tabular}{|c|l|c|l|}
\hline
\multicolumn{2}{|c|}{List-I} & \multicolumn{2}{c|}{List-II}\\[3pt]
\hline
(P) & if $a=1$, $b=2$, then & (1) & $(1-x^2)(1-y^2)=x^2y^2$\\[4pt]
(Q) & if $a=2$, $b=1$, then & (2) & $a+b=8$\\[4pt]
(R) & if $a=0$, $b=1$, then & (3) & \parbox[t]{2.6cm}{$\dfrac{dy}{dx}=-\left(\dfrac{y^2+1}{xy}\right)$, $x,y\in R_0$}\\[10pt]
(S) & \parbox[t]{3.1cm}{if $a=b$ and area bounded by rectangle formed by $f(x,y,a,b)=0$ is 1 sq. unit, then} & (4) & $x^2+4y^2=1+3x^2y^2$\\[3pt]
\hline
\end{tabular}
\end{center}
\begin{qoptions}
\opt{A}{$P\to1;Q\to4;R\to2;S\to1$}
\opt{B}{$P\to1;Q\to4;R\to3;S\to2$}
\opt{C}{$P\to4;Q\to1;R\to3;S\to2$}
\opt{D}{$P\to4;Q\to1;R\to3;S\to1$}
\end{qoptions}
\sect{SECTION-II}
\qn{1} Let $f(x)=\begin{cases}ax\,; & x<2\\ ax^2-bx+3\,; & x\geqslant2\end{cases}$ and `$f$' is differentiable $\forall\,x\in R$, then $(a+b)$ is equal to -
\qn{2} If the real root of the equation $(x-1)^{2019}+(x-2)^{2019}+(x-3)^{2019}+\ldots+(x-2019)^{2019}=0$ is $\lambda$, then $\dfrac{\lambda}{315}$ equals
\qn{3} \[\begin{aligned}f(x)={}&\tan(\sin^{-1}(\cos(\sin^{-1}(x))))\\&\times\tan(\cos^{-1}(\sin(\cos^{-1}x))),\end{aligned}\]where $x\in(0,1)$ then $\displaystyle\sum_{r=0}^{\infty}\dfrac{1}{2^r}f\left(\dfrac{1}{2^r}\right)$ is equal to
\qn{4} The value of\[\lim_{n\to\infty}\left(1-\dfrac1{2^2}\right)\left(1-\dfrac1{3^2}\right)\left(1-\dfrac1{4^2}\right)\ldots\left(1-\dfrac1{n^2}\right)\]is
\qn{5} Let $p_1$ and $p_2$ be two values of $p$ for which $f(x)=\begin{cases}\dfrac{x(\ln(1+x)+\ln(1-x))}{(\sec x-\cos x)}\,; & x\in(-1,0)\\[6pt] \left(p^2-3p-1\right)\sin x+x^2\,; & x\in[0,\infty)\end{cases}$ is differentiable at $x=0$ find $p_1^2+p_2^2$
\qn{6} Let $P(x)$ be a polynomial of degree 4 with $P(1)=1$, $P'(1)=0$, $P''(1)=2$, $P'''(1)=-12$ and $P''''(1)=24$ then value of $\dfrac{P(2)}{P(3)}$ is equal to
\qn{7} A curve is represented by $x=\sec^2t$ and $y=\cot t$, where $t$ is the parameter. If the tangent at point $P$ on curve where $t=\dfrac\pi4$ meets the curve again at $Q$ then $|PQ|=\dfrac{p\sqrt q}{2}$, where $p,q$ are coprime. Find $p+q$
\qn{8} Let $f$ be a real valued function such that $f\left(\dfrac{2x-3}{x-2}\right)=5x-2$, $x\neq2$, then value of $\dfrac54f^{-1}(13)$ is
\end{multicols}
\end{document}
